\documentclass[10pt,a4paper]{amsart}
\usepackage{amsmath,amssymb,amsthm,graphicx}
% Retained preamble for the Stacks Project chapter snapshot in this directory.
% The chapter's own preamble.tex is a copy of the shared Stacks Project preamble; the block
% below is sliced MECHANICALLY, line for line, from it. Nothing was added, removed or
% redefined; the wrapper class, the xr-hyper/hyperref wiring and the externaldocument list
% are omitted because the assembler supplies the document class.
%
% lines taken from preamble.tex: 3, 5, 9, 11, 12, 156, 157, 158, ... (51 lines total)
%
% package coverage restored from the chapter's own preamble (2026-10-13): xy, xy-2cell

\usepackage[all]{xy}
\xyoption{2cell}
\UseAllTwocells

\newenvironment{reference}{\comment}{\endcomment}
\newenvironment{slogan}{\comment}{\endcomment}
\newenvironment{history}{\comment}{\endcomment}
\theoremstyle{plain}
\newtheorem{theorem}[subsection]{Theorem}
\newtheorem{proposition}[subsection]{Proposition}
\newtheorem{lemma}[subsection]{Lemma}
\theoremstyle{definition}
\newtheorem{definition}[subsection]{Definition}
\newtheorem{example}[subsection]{Example}
\newtheorem{exercise}[subsection]{Exercise}
\newtheorem{situation}[subsection]{Situation}
\theoremstyle{remark}
\newtheorem{remark}[subsection]{Remark}
\newtheorem{remarks}[subsection]{Remarks}
\numberwithin{equation}{subsection}
\def\lim{\mathop{\mathrm{lim}}\nolimits}
\def\colim{\mathop{\mathrm{colim}}\nolimits}
\def\Spec{\mathop{\mathrm{Spec}}}
\def\Hom{\mathop{\mathrm{Hom}}\nolimits}
\def\Ext{\mathop{\mathrm{Ext}}\nolimits}
\def\SheafHom{\mathop{\mathcal{H}\!\mathit{om}}\nolimits}
\def\SheafExt{\mathop{\mathcal{E}\!\mathit{xt}}\nolimits}
\def\Sch{\mathit{Sch}}
\def\Mor{\mathop{\mathrm{Mor}}\nolimits}
\def\Ob{\mathop{\mathrm{Ob}}\nolimits}
\def\Sh{\mathop{\mathit{Sh}}\nolimits}
\def\NL{\mathop{N\!L}\nolimits}
\def\CH{\mathop{\mathrm{CH}}\nolimits}
\def\proetale{{pro\text{-}\acute{e}tale}}
\def\etale{{\acute{e}tale}}
\def\QCoh{\mathit{QCoh}}
\def\Ker{\mathop{\mathrm{Ker}}}
\def\Im{\mathop{\mathrm{Im}}}
\def\Coker{\mathop{\mathrm{Coker}}}
\def\Coim{\mathop{\mathrm{Coim}}}
\DeclareMathSymbol{\boxtimes}{\mathbin}{AMSa}{"02}
\def\QCohstack{\mathcal{QC}\!\mathit{oh}}
\def\Cohstack{\mathcal{C}\!\mathit{oh}}
\def\Spacesstack{\mathcal{S}\!\mathit{paces}}
\def\Quotfunctor{\mathrm{Quot}}
\def\Hilbfunctor{\mathrm{Hilb}}
\def\Curvesstack{\mathcal{C}\!\mathit{urves}}
\def\Polarizedstack{\mathcal{P}\!\mathit{olarized}}
\def\Complexesstack{\mathcal{C}\!\mathit{omplexes}}
\def\Pic{\mathop{\mathrm{Pic}}\nolimits}
\def\Picardstack{\mathcal{P}\!\mathit{ic}}
\def\Picardfunctor{\mathrm{Pic}}
\def\Deformationcategory{\mathcal{D}\!\mathit{ef}}
\usepackage{cleveref}
\crefname{originalequation}{original equation}{original equations}
\Crefname{originalequation}{Original equation}{Original equations}
\makeatletter
\newlabel{equation-commutative-square}{{08QM}{source}{}{sourcepointer-1}{}}
\newlabel{equation-commutative-square@cref}{{[originalequation][0][]08QM}{[1][0][]source}{}{}{}}
\makeatother
\crefname{originalequation}{original equation}{original equations}
\Crefname{originalequation}{Original equation}{Original equations}
\makeatletter
\newlabel{equation-double-square}{{08QN}{source}{}{sourcepointer-2}{}}
\newlabel{equation-double-square@cref}{{[originalequation][0][]08QN}{[1][0][]source}{}{}{}}
\makeatother
\crefname{originallemma}{original lemma}{original lemmas}
\Crefname{originallemma}{Original lemma}{Original lemmas}
\makeatletter
\newlabel{lemma-compute-cotangent-complex}{{08PU}{source}{}{sourcepointer-3}{}}
\newlabel{lemma-compute-cotangent-complex@cref}{{[originallemma][0][]08PU}{[1][0][]source}{}{}{}}
\makeatother
\crefname{originallemma}{original lemma}{original lemmas}
\Crefname{originallemma}{Original lemma}{Original lemmas}
\makeatletter
\newlabel{lemma-flat-base-change}{{08QP}{source}{}{sourcepointer-4}{}}
\newlabel{lemma-flat-base-change@cref}{{[originallemma][0][]08QP}{[1][0][]source}{}{}{}}
\makeatother
\crefname{originallemma}{original lemma}{original lemmas}
\Crefname{originallemma}{Original lemma}{Original lemmas}
\makeatletter
\newlabel{sites-cohomology-lemma-change-of-rings}{{08Q8}{source}{}{sourcepointer-5}{}}
\newlabel{sites-cohomology-lemma-change-of-rings@cref}{{[originallemma][0][]08Q8}{[1][0][]source}{}{}{}}
\makeatother
\title{Flat base change for the cotangent complex}
\author{The Stacks Project Authors}
\date{stacks-project:cotangent.tex:local-snapshot-20261012; GNU FDL 1.2 or later}
\begin{document}
\maketitle
\setcounter{subsection}{1}
\noindent\textit{Editorial note.} All delivered mathematical text is the author's own, copied from the pinned Stacks Project chapter file cotangent.tex; no statement or proof step was rewritten, reordered or removed.\par\medskip
\noindent\textit{Editorial note.} The statement refers to the commutative square and the proof to the double square of the chapter; both are fixed by labelled equations and the displayed computation is inside the counted span.\par\medskip
\noindent\textit{Editorial note.} The five references of the counted spans are recorded as non-load-bearing navigation pointers; the argument that consumes them stays inside the counted proof span.\par\medskip
\noindent\textit{Editorial note.} Preamble: the retained preamble of this chapter (sources/stacks-cotangent/retained-preamble.tex, lines 1-63), containing the Xy-pic setup, the comment environments, the theorem declarations and the macros including \textbackslash{}NL.\par\medskip
\noindent\textit{Allowed standard prerequisites (editorial).}\begin{itemize}
\item The computation of the cotangent complex as a lower shriek of differentials. Lemma lemma-compute-cotangent-complex (545-555, produced as L44157), used at proof lines 1215-1220 and 1232-1233. Rewrites both sides as lower shrieks of differential sheaves.
\item The change-of-rings lemma for lower shrieks. Cohomology on Sites, Lemma sites-cohomology-lemma-change-of-rings, used at proof lines 1221-1222. Turns the derived tensor product into a lower shriek along h.
\item Flatness of the sheaf of differentials over the structure sheaf. used at proof lines 1223-1230. Computes the pullback of the differential sheaf as the differential sheaf of the primed square.
\item Flat base change for the lower shriek in this setting. Lemma lemma-flat-base-change (1174-1181, produced as L44146), used at proof lines 1231-1233. Identifies the two lower shrieks and hence concludes the isomorphism.
\end{itemize}\medskip
\ifdefined\hypertarget\hypertarget{sourcepointer-1}{}\fi
\noindent\textit{Source reference.} Original equation 08QM: focused-seeds/AGP/sources/stacks-snapshot-20261011/tags.txt, source lines 11065--11065. Navigation pointer into the Stacks Project; the standard fact it supplies is applied inside the counted proof span and the argument resting on it stays there.\par\medskip
\ifdefined\hypertarget\hypertarget{sourcepointer-2}{}\fi
\noindent\textit{Source reference.} Original equation 08QN: focused-seeds/AGP/sources/stacks-snapshot-20261011/tags.txt, source lines 11066--11066. Navigation pointer into the Stacks Project; the standard fact it supplies is applied inside the counted proof span and the argument resting on it stays there.\par\medskip
\ifdefined\hypertarget\hypertarget{sourcepointer-3}{}\fi
\noindent\textit{Source reference.} Original lemma 08PU: focused-seeds/AGP/sources/stacks-snapshot-20261011/tags.txt, source lines 11037--11037. Navigation pointer into the Stacks Project; the standard fact it supplies is applied inside the counted proof span and the argument resting on it stays there.\par\medskip
\ifdefined\hypertarget\hypertarget{sourcepointer-4}{}\fi
\noindent\textit{Source reference.} Original lemma 08QP: focused-seeds/AGP/sources/stacks-snapshot-20261011/tags.txt, source lines 11067--11067. Navigation pointer into the Stacks Project; the standard fact it supplies is applied inside the counted proof span and the argument resting on it stays there.\par\medskip
\ifdefined\hypertarget\hypertarget{sourcepointer-5}{}\fi
\noindent\textit{Source reference.} Original lemma 08Q8: focused-seeds/AGP/sources/stacks-snapshot-20261011/tags.txt, source lines 11051--11051. Navigation pointer into the Stacks Project; the standard fact it supplies is applied inside the counted proof span and the argument resting on it stays there.\par\medskip
\begin{lemma}
\label{lemma-flat-base-change-cotangent-complex}
If (\ref{equation-commutative-square}) induces a quasi-isomorphism
$B \otimes_A^\mathbf{L} A' = B'$, then the functoriality map
induces an isomorphism
$$
L_{B/A} \otimes_B^\mathbf{L} B' \longrightarrow L_{B'/A'}
$$
\end{lemma}

\begin{proof}
We will use the notation introduced in Equation (\ref{equation-double-square}).
We have
$$
L_{B/A} \otimes_B^\mathbf{L} B' =
L\pi_!(\Omega_{\mathcal{O}/A} \otimes_\mathcal{O} \underline{B})
\otimes_B^\mathbf{L} B' =
L\pi_!(Lh^*(\Omega_{\mathcal{O}/A} \otimes_\mathcal{O} \underline{B}))
$$
the first equality by Lemma \ref{lemma-compute-cotangent-complex}
and the second by Cohomology on Sites, Lemma
\ref{sites-cohomology-lemma-change-of-rings}.
Since $\Omega_{\mathcal{O}/A}$ is a flat $\mathcal{O}$-module,
we see that $\Omega_{\mathcal{O}/A} \otimes_\mathcal{O} \underline{B}$
is a flat $\underline{B}$-module. Thus
$Lh^*(\Omega_{\mathcal{O}/A} \otimes_\mathcal{O} \underline{B}) =
\Omega_{\mathcal{O}/A} \otimes_\mathcal{O} \underline{B'}$
which is equal to
$g^{-1}(\Omega_{\mathcal{O}'/A'} \otimes_{\mathcal{O}'} \underline{B'})$
by inspection.
we conclude by Lemma \ref{lemma-flat-base-change}
and the fact that $L_{B'/A'}$ is computed by
$L\pi'_!(\Omega_{\mathcal{O}'/A'} \otimes_{\mathcal{O}'} \underline{B'})$.
\end{proof}

\end{document}
